% !TeX root = ../main.tex

\section{Ethical Considerations}
\label{sec:ethics}

\paragraph{Exposure to harmful content.}
The datasets used in this study contain memes that may include
hateful, discriminatory, misogynistic, offensive, politically
sensitive, or otherwise disturbing content.
Researchers who inspect raw memes, retrieved evidence, generated
rationales, or error cases may therefore be exposed to potentially
harmful material.
We minimize unnecessary exposure by separating automated
experiments from manual review, limiting qualitative inspection to
samples required for validation or analysis, and avoiding the
inclusion of gratuitous harmful examples in the paper.
Any future human evaluation should provide annotators with an
explicit content warning, allow them to pause or withdraw, and
avoid requiring repeated exposure to particularly severe content.

\paragraph{Dataset provenance and privacy.}
We use existing meme benchmarks according to their intended
research purposes and retain their original sample identifiers,
dataset families, and domain metadata.
The source protocol combines the COVID-19 and U.S.\ politics
domains of HarMeme for training and validation, while FHM is used
only as a held-out target dataset.
We do not use FHM samples for training, prompt development,
threshold selection, retrieval-index construction, or any other
model-selection decision.

Memes may contain identifiable people, public figures, usernames,
logos, or contextual references.
Our task is content-level harmfulness interpretation rather than
person identification.
The framework does not attempt to infer private identity,
demographic attributes, or sensitive personal information beyond
what is explicitly required by the dataset annotation schema.
Any public release of code or experiment manifests should avoid
redistributing raw images unless permitted by the licenses and
terms of the original datasets.

% TODO_ETHICS: Verify the final redistribution and licensing policy
% of every dataset before releasing code or artifacts.

\paragraph{Annotation provenance and uncertainty.}
The harmfulness labels used in this study originate from the
respective source benchmarks.
The target, intent, and tactic annotations are produced through the
existing structured annotation pipeline and are preserved with
their confidence, audit flags, masks, and provenance metadata.
In particular, the current structured annotations for FHM are
agent-generated, schema-validated silver annotations rather than
human-adjudicated gold labels.
We therefore report the corresponding results as
\emph{silver structured evaluation} and provide field-level
coverage and valid-sample counts.

Automated annotation may reproduce the cultural assumptions,
social biases, and systematic errors of the underlying models.
Labels such as harmfulness, intent, and rhetorical tactic may also
admit multiple reasonable interpretations.
We retain explicit \texttt{ambiguous}, \texttt{unknown}, and
masked states where defined by the annotation protocol, rather than
forcing every sample into a confident category.
Human-evaluation results are not reported unless independently
collected annotations have been validated and imported through the
designated evaluation procedure.

\paragraph{External knowledge and contextual bias.}
External contextual information can help interpret implicit
references, but it may also be incomplete, outdated, culturally
specific, or factually incorrect.
Retrieved passages, generated hypotheses, and fallback candidates
are therefore stored as distinct candidate types and are not
treated as verified knowledge by default.
The Stage C verifier filters candidates according to relevance,
validity, and task-specific support; however, this procedure does
not constitute a guarantee of factual correctness.

Knowledge sources may disproportionately represent dominant
languages, regions, institutions, or political perspectives.
Consequently, a system that relies on external retrieval may
perform unevenly across communities or cultural contexts.
We retain source and candidate provenance to make such failures
auditable and include cultural-context errors in the planned error
analysis.

\paragraph{Rationale limitations.}
Generated rationales are intended to make model decisions easier
to inspect, not to provide authoritative factual explanations.
A rationale may be fluent while still omitting relevant context,
overstating weak evidence, or introducing an unsupported claim.
For this reason, rationale text is not used as the formal source of
target, intent, or tactic predictions.
Formal structured metrics are computed from trainable model
outputs, while rationale quality is evaluated separately through
grounding, support, and unsupported-claim analyses.

The evidence attribution and attention traces retained by the
framework should also not be interpreted as causal explanations.
They are diagnostic records that indicate which evidence was
selected or emphasized by the implemented pipeline.

\paragraph{Risks of automated moderation.}
False positives may incorrectly flag satire, criticism, political
speech, or benign humor, whereas false negatives may allow harmful
content to remain undetected.
Errors may be especially consequential for marginalized groups,
whose language and cultural references are often underrepresented
in training data.
Accordingly, the proposed framework should not be deployed as an
autonomous moderation authority.
Its predictions, evidence records, and rationales are better viewed
as decision-support information for trained human moderators.

The structured outputs should likewise not be used to assign intent
or harmful attributes to individuals outside the context of the
analyzed meme.
Model confidence does not remove the need for human review,
particularly in ambiguous, culturally dependent, or high-impact
cases.

\paragraph{Dual-use considerations.}
The same analysis capabilities that support moderation research
could be misused to identify weaknesses in content-moderation
systems or to create memes designed to evade detection.
To reduce this risk, released artifacts should prioritize model
interfaces, evaluation protocols, and aggregated results rather than
unnecessary reproduction of harmful material.
We also avoid presenting retrieved or generated harmful content as
endorsed factual information.

\paragraph{Human evaluation and oversight.}
The current automated pipeline can export annotation packages for
evidence relevance, rationale quality, verifier decisions, and
stage-level error analysis.
It does not generate human gold judgments automatically.
Any future human evaluation should use informed participation,
clear annotation guidelines, appropriate compensation, secure data
handling, and a process for resolving annotator disagreements.
Institutional review or ethics approval requirements should be
determined before such evaluation is conducted.

% TODO_ETHICS: Add the final human-evaluation protocol, annotator
% demographics at an aggregate level, compensation policy, and
% institutional review status if human evaluation is performed.